%%
% BIThesis 研究生学位论文模板 The BIThesis Template for Graduate Thesis
% This file has no copyright assigned and is placed in the Public Domain.
%%

\begin{abstract}
  为满足海洋信息网络和空天地海一体化体系中日益增长的水下信息传输需求，本文围绕水下无线光通信在低动态浑浊多径场景和高动态移动场景中的可靠传输问题开展研究。针对不同传播条件下主导失真机制差异显著、传统正交频分复用方案适应性不足的问题，本文以仿射频分复用为统一技术主线，对波形构造、参数设计和检测方法进行了系统研究。本文主要工作与结果如下：

  （1）针对低动态浑浊多径场景中传统光正交频分复用在非负实信号约束下对时延扩展和频率选择性衰落适应性不足的问题，提出面向强度调制直接检测架构的非对称限幅光仿射频分复用方案。该方案改善了传统方案对多径传播和深衰落位置的敏感性，提升了误码率性能与多径鲁棒性；仿真结果表明，在双径条件下相较非对称限幅光正交频分复用约获得 2 分贝增益，在路径数由 2 增至 10 时性能退化更缓。

  （2）针对高动态移动场景中多普勒扩展和相位噪声共同作用下传统方案频偏失配严重的问题，提出面向相干接收的仿射频分复用索引调制方案，并利用静默位置构造残余频偏补偿方法。该方法改善了高动态条件下主导频偏难以低开销补偿的问题，提升了系统在多普勒扩展和激光器线宽扰动下的检测可靠性；仿真结果表明，所提方案在归一化多普勒频移 0 至 0.20 的范围内保持可靠通信，并在高动态条件下优于传统索引调制方案。

  （3）与传统显式导频补偿方案相比，所提方法无需额外导频开销，改善了有效频谱利用率和能量分配效率，避免了 20\% 导频占用带来的约 0.97 分贝等效信噪比损失。
  总体而言，本文形成了面向低动态与高动态两类典型场景的统一设计主线，可为高可靠水下光通信系统设计提供理论参考。 
\end{abstract}

% 如需手动控制换行连字符位置，可写 aa\-bb，详见
% https://bithesis.bitnp.net/faq/hyphen.html

\begin{abstractEn}
  To meet the growing demand for underwater information transmission in marine information networks and integrated space-air-ground-sea systems, this thesis studies reliable transmission for low-dynamic turbid multipath scenarios and high-dynamic mobile scenarios in underwater optical communication. To address the limited adaptability of conventional orthogonal frequency division multiplexing under different dominant impairments, affine frequency division multiplexing is adopted as the unified technical line, and waveform construction, parameter design, and detection methods are investigated systematically. The main improvements and conclusions are as follows:

  (1) For low-dynamic turbid multipath scenarios, an asymmetrically clipped optical affine frequency division multiplexing scheme is proposed for intensity-modulation direct-detection links. The proposed scheme reduces the sensitivity of conventional optical orthogonal frequency division multiplexing to delay spread, frequency-selective fading, and the nonnegative real-signal constraint, thereby improving bit-error-rate performance and multipath robustness. Simulation results show a gain of about 2 decibels over asymmetrically clipped optical orthogonal frequency division multiplexing under the two-path condition, with slower performance degradation when the number of paths increases from 2 to 10.

  (2) For high-dynamic mobile scenarios, a coherent transmission scheme based on affine frequency division multiplexing and index modulation is proposed, together with a residual carrier-frequency-offset compensation method constructed from silent positions. The proposed method alleviates severe frequency-offset mismatch under joint Doppler spread and phase-noise disturbances, and improves detection reliability under Doppler variation and laser-linewidth perturbations. Simulation results show reliable communication over normalized Doppler shifts from 0 to 0.20, while achieving better performance than the conventional index-modulation benchmark in high-dynamic conditions.

  (3) Compared with the conventional pilot-assisted compensation strategy, the proposed method requires no additional pilot overhead, improves effective spectral utilization and power allocation efficiency, and avoids the equivalent signal-to-noise-ratio loss of about 0.97 decibels caused by 20\% pilot occupation. Overall, this thesis establishes a unified scenario-oriented design route for representative low-dynamic and high-dynamic underwater links, and provides theoretical support for reliable underwater optical communication system design.
\end{abstractEn}
